\clearpage\thispagestyle{empty}
\pdfbookmark[0]{Sources for the report on Lindenau’s Mercury tables}{lindenaumercury1813-sources}
\begin{center}\Large\bfseries Sources for the report on Lindenau’s Mercury tables\end{center}

The translation follows Carl Friedrich Gauss’s German report in \emph{Werke}, vol.~VI (Göttingen, 1874), pp.~559--560. It first appeared in \emph{Göttingische gelehrte Anzeigen}, no.~107, 5 July 1813, pp.~1065--1069. The Göttingen State and University Library holds the digitized journal.

In the collected text as transcribed here, the denominator of the fraction of the Sun’s mass is 849132. The journal printing on p.~1068 gives 349132. The report quotes the title of a forthcoming work. The printed title and the results in Bernhard von Lindenau’s \emph{Investigatio nova orbitae a Mercurio circa Solem descriptae} (Gotha, Libraria Beckeriana, 1813) differ in some particulars. On p.~30 Lindenau prints the factor $1{,}09643$ and the fraction $\frac{1}{349440}$, and p.~31 gives his own table of elements.

AI-translated from German.

First printing (\href{https://images.p.sub.uni-goettingen.de/iiif/image/gdz:PPN319721507_1813:00001069/full/1650,/0/default.jpg}{1065}, \href{https://images.p.sub.uni-goettingen.de/iiif/image/gdz:PPN319721507_1813:00001070/full/1650,/0/default.jpg}{1066}, \href{https://images.p.sub.uni-goettingen.de/iiif/image/gdz:PPN319721507_1813:00001071/full/1650,/0/default.jpg}{1067}, \href{https://images.p.sub.uni-goettingen.de/iiif/image/gdz:PPN319721507_1813:00001072/full/1650,/0/default.jpg}{1068}, \href{https://images.p.sub.uni-goettingen.de/iiif/image/gdz:PPN319721507_1813:00001073/full/1650,/0/default.jpg}{1069}).

Lindenau’s book (\href{https://doi.org/10.3931/e-rara-72976}{ETH-Bibliothek Zürich, Rar 33093}, \href{https://www.e-rara.ch/i3f/v20/20184929/full/1650,/0/default.jpg}{title page}, \href{https://www.e-rara.ch/i3f/v20/20184958/full/1650,/0/default.jpg}{30}, \href{https://www.e-rara.ch/i3f/v20/20184959/full/1650,/0/default.jpg}{31}).

\clearpage
